\documentclass[a4paper,11pt]{amsart}

\usepackage[margin=1in]{geometry}
\usepackage{amsmath,amssymb,amsthm}

\title{Distinct Dyadic Subsum Representations}
\author{}
\date{}

\newtheorem{theorem}{Theorem}
\newtheorem{lemma}{Lemma}
\newtheorem{proposition}{Proposition}
\newtheorem{corollary}{Corollary}
\newtheorem{remark}{Remark}

\begin{document}

\maketitle

\section{The finite problem}

Put
\[
        u_m=\frac{m}{2^m}.
\]
We first prove, unconditionally, that infinitely many positive integers \(n\)
admit a representation
\[
        u_n=\sum_{k=1}^t u_{a_k},
        \qquad t\geq 2,
\]
with \(a_1,\ldots,a_t\) distinct positive integers.

\begin{lemma}
\label{monotone}
For \(m\geq 2\), the sequence \(m/2^m\) is strictly decreasing.
Consequently, if \(n\geq 3\) and
\[
        \frac n{2^n}=\sum_{a\in A}\frac a{2^a}
\]
with \(A\) finite and all summands positive, then every \(a\in A\) satisfies
\(a>n\).
\end{lemma}

\begin{proof}
The inequality
\[
        \frac{m+1}{2^{m+1}}<\frac m{2^m}
\]
is equivalent to \(m+1<2m\), which holds for \(m\geq 2\).
Thus for \(n\geq 3\), any summand with \(a<n\) would already exceed
\(n/2^n\), while a summand with \(a=n\) would exhaust the whole left side and
leave no room for the remaining positive summands.
\end{proof}

\begin{theorem}
\label{infinite-family}
For every integer \(r\geq 2\), let
\[
        n_r=2^{r+1}-r-2.
\]
Then
\[
        \frac{n_r}{2^{n_r}}
        =
        \sum_{j=1}^{r}\frac{n_r+j}{2^{n_r+j}}.
\]
In particular, infinitely many \(n\) have the required finite distinct
representation.
\end{theorem}

\begin{proof}
For any positive integer \(n\),
\[
        \sum_{j=1}^{r}\frac{n+j}{2^{n+j}}
        =
        \frac1{2^n}\sum_{j=1}^{r}\frac{n+j}{2^j}.
\]
Now
\[
        \sum_{j=1}^{r}\frac1{2^j}=1-\frac1{2^r}
\]
and
\[
        \sum_{j=1}^{r}\frac j{2^j}
        =
        2-\frac{r+2}{2^r}.
\]
Therefore
\[
        \sum_{j=1}^{r}\frac{n+j}{2^j}
        =
        n\left(1-\frac1{2^r}\right)
        +
        2-\frac{r+2}{2^r}.
\]
This equals \(n\) precisely when
\[
        n=2^{r+1}-r-2.
\]
Thus, for \(n=n_r\),
\[
        \frac{n_r}{2^{n_r}}
        =
        \sum_{j=1}^{r}\frac{n_r+j}{2^{n_r+j}}.
\]
The indices \(n_r+1,\ldots,n_r+r\) are distinct, and the integers \(n_r\) are
strictly increasing for \(r\geq 2\).
\end{proof}

The first cases are
\[
        \frac4{2^4}
        =
        \frac5{2^5}+\frac6{2^6},
\]
and
\[
        \frac{11}{2^{11}}
        =
        \frac{12}{2^{12}}+\frac{13}{2^{13}}+\frac{14}{2^{14}}.
\]

\section{The all-\(n\) question}

The construction above proves infinitely many cases.  It does not prove that
every \(n\) works.  Here is the exact finite reformulation.

\begin{proposition}
\label{subset-condition}
Let \(n\geq 3\).  Then \(u_n\) has a finite distinct representation by at least
two other terms if and only if there is a finite set
\[
        J\subseteq \{1,2,\ldots\},\qquad |J|\geq 2,
\]
such that
\[
        n=\sum_{j\in J}\frac{n+j}{2^j}.
\]
Equivalently,
\[
        n\left(1-\sum_{j\in J}2^{-j}\right)
        =
        \sum_{j\in J}j2^{-j}.
\]
\end{proposition}

\begin{proof}
By Lemma \ref{monotone}, every summand in a representation of \(u_n\) has the
form \(u_{n+j}\) with \(j\geq 1\).  Thus
\[
        \frac n{2^n}
        =
        \sum_{j\in J}\frac{n+j}{2^{n+j}}
\]
for a finite set \(J\).  Multiplying by \(2^n\) gives
\[
        n=\sum_{j\in J}\frac{n+j}{2^j}.
\]
Separating the terms involving \(n\) gives the displayed equivalent identity.
The converse is obtained by reversing the same calculation.
\end{proof}

It is useful to spell out the residual graph that would settle the all-\(n\)
question.  Fix \(n\geq 2\).  Suppose decisions have been made for offsets
\(1,\ldots,j\), and let \(J_j\subseteq \{1,\ldots,j\}\) be the selected
offsets.  Define
\[
        y_{j+1}
        =
        2^j\left(
        n-\sum_{i\in J_j}\frac{n+i}{2^i}
        \right).
\]
Then \(y_1=n\).  At offset \(j\), omitting \(n+j\) gives
\[
        y_{j+1}=2y_j,
\]
while including \(n+j\) gives
\[
        y_{j+1}=2y_j-(n+j).
\]
The invariant
\[
        0\leq y_j<n+j
\]
is preserved if at each step \(y_{j+1}\) is chosen as the least nonnegative
residue of \(2y_j\) modulo \(n+j\).  Indeed,
\[
        0\leq 2y_j<2(n+j),
\]
so this residue is one of the two allowed moves above.  If some \(y_{j+1}=0\),
then the selected offsets give the identity in Proposition
\ref{subset-condition}.

Writing
\[
        h_j=y_j,\qquad d_j=n+j-y_j,
\]
one obtains \(h_j+d_j=n+j\).  If \(h_j>d_j\), then the next move is
\[
        (h_j,d_j)\mapsto (h_j-d_j,\,2d_j+1),
\]
whereas if \(h_j<d_j\), the next move is
\[
        (h_j,d_j)\mapsto (2h_j,\,d_j-h_j+1).
\]
If \(h_j=d_j\), then \(2y_j=n+j\), and including \(n+j\) makes the next
residual zero.

\begin{remark}
Thus the assertion that every \(n\geq 2\) works would follow from the
termination statement that the above recurrence, started at \((h_1,d_1)=(n,1)\),
eventually reaches a diagonal pair.  This note does not prove that termination
statement, so no claim that all \(n\) work is made here.
\end{remark}

\section{Infinite subsums}

We now prove the requested infinite-representation statement independently of
the finite all-\(n\) question.

\begin{lemma}
\label{tail}
For \(N\geq 1\),
\[
        \sum_{m=N}^{\infty}\frac m{2^m}
        =
        \frac{N+1}{2^{N-1}}.
\]
Consequently,
\[
        \sum_{m=a+1}^{\infty}\frac m{2^m}
        =
        \frac{a+2}{2^a}
        >
        \frac a{2^a}
\]
for every \(a\geq 1\).
\end{lemma}

\begin{proof}
The identity follows from the usual differentiated geometric series, or by
checking that the right side has the same successive differences as the tail:
\[
        \frac{N+1}{2^{N-1}}-\frac{N+2}{2^N}
        =
        \frac N{2^N}.
\]
The inequality is immediate from \(a+2>a\).
\end{proof}

\begin{lemma}[Overlap subsum theorem]
\label{overlap}
Let \((b_m)\) be a positive summable sequence with total \(S\), and put
\[
        T_m=\sum_{k=m}^{\infty}b_k.
\]
Assume that
\[
        b_m<T_{m+1}
\]
for every \(m\).  Then every point of the open interval \((0,S)\) has at least
\(2^{\aleph_0}\) representations as a subsum
\[
        x=\sum_{m\in A}b_m.
\]
\end{lemma}

\begin{proof}
We recall the standard nested-interval argument.  After a finite set of
decisions before index \(m\), the possible remaining values form an interval of
length \(T_m\).  At index \(m\), the two choices give the intervals
\[
        [0,T_{m+1}]
        \qquad\text{and}\qquad
        [b_m,b_m+T_{m+1}],
\]
translated by the already chosen partial sum.  These intervals overlap in a
nondegenerate interval because \(b_m<T_{m+1}\).

Fix \(x\in(0,S)\).  Starting from the full interval \([0,S]\), recursively
choose deeper and deeper finite cylinders containing \(x\).  Since every tail
interval has nondegenerate overlap at its next split and the tail lengths tend
to zero, each cylinder containing \(x\) has a later descendant at which \(x\)
lies in the overlap of the two child cylinders.  Keeping both children at each
such branching stage embeds the full binary tree into the set of subsum
representations of \(x\).  Distinct infinite paths through this tree differ at
the first branch where they choose different children, hence give distinct
subsets.  The set of infinite binary paths has cardinality \(2^{\aleph_0}\).
\end{proof}

\begin{theorem}
There exists a rational number \(x\) with at least \(2^{\aleph_0}\) distinct
representations
\[
        x=\sum_{k=1}^{\infty}\frac{a_k}{2^{a_k}},
\]
where the \(a_k\) are distinct positive integers.  For example, one may take
\[
        x=1.
\]
\end{theorem}

\begin{proof}
Apply Lemma \ref{overlap} to
\[
        b_m=\frac m{2^m}.
\]
By Lemma \ref{tail}, each term is strictly smaller than the remaining tail, and
the total sum is
\[
        \sum_{m=1}^{\infty}\frac m{2^m}=2.
\]
Therefore every \(x\in(0,2)\), in particular the rational number \(x=1\), has
at least continuum-many subsum representations.  Listing the selected set
\(A\subseteq \mathbb N\) increasingly as \(A=\{a_1<a_2<\cdots\}\) gives the
required form.
\end{proof}

\end{document}
